\documentclass[12pt]{article}
\usepackage[utf8]{inputenc}
\usepackage[T1]{fontenc}
\usepackage[margin=1in]{geometry}
\usepackage{setspace}
\usepackage{xcolor}
\usepackage[colorlinks=true, urlcolor=blue!60!black, linkcolor=blue!60!black]{hyperref}
\setlength{\parindent}{0em}
\setlength{\parskip}{0.6em}
\renewcommand{\baselinestretch}{1.0}
\begin{document}
\begin{center}
{\LARGE \textit{Board of Trustees of the University of Alabama v. Garrett}, 531 U.S. 356 (2001)\par}
\vspace{0.5em}
{\large Case Brief --- Two-Page Edition\par}
\end{center}

\section{Memory Jogger}
State employees cannot sue their state employers for money damages under Title I of the ADA: Congress showed no pattern of unconstitutional state discrimination, and the remedy was not ``congruent and proportional'' (\href{https://www.courtlistener.com/opinion/118407/board-of-trustees-of-univ-of-ala-v-garrett/}{CourtListener}).

\section{Detailed Case Facts}
Nurse Patricia Garrett was told to give up her Director of Nursing post at the University of Alabama in Birmingham Hospital after breast-cancer treatment, and she transferred to a lower-paying job. Security officer Milton Ash was denied accommodations for asthma and sleep apnea, and his evaluations dropped after he filed an EEOC charge. Title I requires ``reasonable accommodations'' and expressly seeks to abrogate state immunity.

\section{Procedural History}
The Northern District of Alabama granted summary judgment for the State; the Eleventh Circuit reversed. The Supreme Court reversed as to Title I and dismissed the Title II question as improvidently granted.

\section{Judicial Votes}
Reversed 5--4. Chief Justice Rehnquist wrote for O'Connor, Scalia, Kennedy, and Thomas; Kennedy, joined by O'Connor, concurred; Breyer, joined by Stevens, Souter, and Ginsburg, dissented.

\section{Holding}
The Eleventh Amendment bars private damages suits against States under Title I. For rights subject only to rational-basis review, Congress must document a pattern of unconstitutional discrimination by the States themselves and enact congruent and proportional remedies. Title I's standards still bind States through federal enforcement and \textit{Ex parte Young} injunctions.

\section{Analysis of Opinions}
\textbf{Majority:} Under \textit{Cleburne}, disability classifications need only be rational. Excluding local governments, which lack immunity, the record held only ``half a dozen examples'' and unexamined anecdotes, not a pattern of state violations, and Title I's accommodation and disparate-impact rules far exceed equal protection. \textbf{Kennedy:} the ADA is a moral milestone, but charging sovereign States with purposeful wrongdoing requires ``confirming judicial documentation.'' \textbf{Breyer:} ``a legislature is not a court of law''; rational-basis deference restrains courts, not Congress, and Section 5 permits any rational means.

\section{Examples: Future Applications}
\textbf{Same Side:} damages suits for genetic-status discrimination, on a record of private employers and four state anecdotes, are barred. \textbf{Opposite Side:} courthouse-access suits under Title II survive because they enforce a fundamental right (\textit{Tennessee v. Lane}). \textbf{Fence-Sitter:} damages suits over inaccessible state-university websites, resting on 60 targeted complaints and a narrower remedy.

\section{Critique}
Critics call \textit{Garrett} judicial activism that turns Section 5 into a record test a 1990 Congress could not anticipate. Textualists note the Eleventh Amendment does not cover suits by a State's own citizens, and fusing \textit{Cleburne} with \textit{Boerne} leaves Congress almost unable to protect non-suspect classes.

\section{Key Quotations}
\begin{quote}
``Thus, the result of \textit{Cleburne} is that States are not required by the Fourteenth Amendment to make special accommodations for the disabled, so long as their actions toward such individuals are rational. They could quite hardheadedly---and perhaps hardheartedly---hold to job-qualification requirements which do not make allowance for the disabled. If special accommodations for the disabled are to be required, they have to come from positive law and not through the Equal Protection Clause.''
\end{quote}
\noindent --- Chief Justice Rehnquist, Opinion of the Court, 531 U.S. at 367--368.

\end{document}
